%% kapa: exam=1 round=22 year=2011 period=2 form=A subject=appraisal_law subject_ko=부동산관계법규 seq=9 number=49
%% source: Q-Net 감정평가사 기출문제 (cst00309) artlSeq=5036790 file=extracted/1008/5036790/2011년 제22회 감정평가사 1차시험 기출문제/2교시.hwp
%% conversion: kordoc 4.24.51 -> tools/gen1.py (auto); review=proofread-gpt-astra; flags=none
\begin{question}{49}
국토의 계획 및 이용에 관한 법령상 지구단위계획구역으로 지정하려는 구역에서 도시관리계획 입안시 토지적성평가가 면제될 수 있는 경우가 아닌 것은?
\begin{choices}
  \item 도시관리계획의 입안 내용이 개발용도의 용도지역 상호간의 변경인 경우
  \item 당해 지구단위계획구역안의 나대지면적이 구역면적의 2퍼센트에 미달하는 경우
  \item 당해 지구단위계획구역이 상업지역과 상업지역에 연접한 지역에 위치하는 경우
  \item 주거지역ㆍ상업지역 또는 공업지역에 도시관리계획을 입안하는 경우
  \item 도시계획시설부지에서 도시관리계획을 입안하는 경우
\end{choices}
\end{question}
\answer{1}
